\section{Fundamentals and State of Research} \label{sec:fundamentals}

This section provides an overview of the fundamental concepts and the current state of research relevant to this thesis. It covers key topics in explainable artificial intelligence (XAI) in section \ref{ssec:xai} and cognitive psychology in section \ref{ssec:cognitive-psychology}. Related works are discussed in section \ref{ssec:related-works}, and the research objectives of this thesis are outlined in section \ref{ssec:research-objectives}.

\subsection{Explainable Artificial Intelligence} \label{ssec:xai}



\subsubsection{Terminology} \label{sssec:terminology}

In this thesis, \enquote{explainability} is used as the preferred term, while \enquote{interpretability} is treated as the broader umbrella concept. This decision follows recent \ac{XAI} literature, where both terms are sometimes used interchangeably, but are conceptually distinct \parencite{Arrieta2020,Guidotti2018}. Building on \textcite{Arrieta2020}, an \ac{AI} system is explainable if it can provide details or reasons that make its functioning understandable for a specific audience. The audience perspective is essential: what is understandable for an expert user may be insufficient for a lay user \parencite{Arrieta2020,Lipton2018}.

Interpretability refers to the extent to which a human can assign meaning to a model, its behaviour, or its output \parencite{Arrieta2020}. Explainability, in contrast, focuses on the explicit communication layer between system and user, where reasons are provided in a form that supports human understanding and decision-making \parencite{Guidotti2018,Lipton2018}. Transparency is used here as a complementary concept describing how directly a model can be inspected without additional explanation interfaces. Following \textcite{Lipton2018}, transparency can be discussed at the level of the overall model, individual components, and the training procedure.

Insights from cognitive psychology sharpen this terminology. Explanations are typically framed as answers to why-questions and are evaluated by whether they satisfy the explainee's information need \parencite{Miller2019,Overton2011}. In practice, people often prefer concise, contrastive explanations over complete causal chains, because full chains are both cognitively and computationally expensive \parencite{Lipton1990,Kelley1987}. This is relevant for \ac{XAI}, where explanation quality is not only technical correctness, but also perceived usefulness for the target audience \parencite{Miller2019,Lombrozo2006,Kashima1998}.

\begin{figure}[ht]
    \centering
    % TODO: Insert timeline figure showing publication trends for terminology keywords.
    \fbox{\parbox[c][4cm][c]{0.85\linewidth}{\centering Placeholder for terminology timeline figure}}
    \caption{Placeholder timeline for the usage of terminology in \ac{XAI} literature.}
    \label{fig:terminology-timeline-placeholder}
\end{figure}

Table~\ref{tab:question-types-explanations} summarises three common explanation question types used throughout this thesis.

\begin{ctable}
    \centering
    \begin{tabularx}{\textwidth}{l|X|X}
        \textbf{Question} & \textbf{Reasoning} & \textbf{Description} \\
        \hline
        What & Associative & Reasons about which unobserved events may have occurred given the observed outcome. \\
        How & Interventionist & Simulates a change in the situation to assess whether the outcome would still occur. \\
        Why & Counterfactual & Compares alternative causes to explain why one outcome happened instead of another.
    \end{tabularx}
    \caption[Types of Explanations]{Types of explanations as a product and associated reasoning, adapted from \textcite{Miller2019}.}
    \label{tab:question-types-explanations}
\end{ctable}

Compact vignette used in this thesis: a student receives an incorrect recommendation from a learning-support model. A \enquote{what} explanation lists the features associated with that recommendation, a \enquote{how} explanation shows what changes in input would alter the result, and a \enquote{why} explanation contrasts the actual recommendation with the most plausible alternative.

\subsubsection{Generative vs. Discriminative Models} \label{sssec:generative-vs-discriminative}

\subsubsection{Explainability Methods} \label{sssec:xai-methods}

\subsection{Cognitive Psychology} \label{ssec:cognitive-psychology}

\subsubsection{Dual Process Theory} \label{sssec:dual-process-theory}

\subsubsection{Heuristics} \label{sssec:heuristics}

\subsubsection{Theory of Mind} \label{ssec:theory-of-mind}

\subsubsection{Theory of Mind in \acs{XAI}} \label{ssec:theory-of-mind-xai}

\subsection{Related Works} \label{ssec:related-works}

\subsection{Research Objectives} \label{ssec:research-objectives}

\subsection{Filler}

The following sources (among others) will be used for the fundamentals section.

\begin{itemize}
    \item \textcite{Arrieta2020}
    \item \textcite{Lipton2018}
    \item \textcite{Lipton1990}
    \item \textcite{Lundberg2017}
    \item \textcite{Miller2019}
    \item \textcite{Ribeiro2016}
    \item \textcite{Schneider2024}
    \item \textcite{Wachter2017}
\end{itemize}
